\documentclass[11pt,a4paper]{article}
\usepackage{fontspec,amsmath,amssymb,amsthm,microtype}
\setmainfont{lmroman10-regular.otf}[BoldFont=lmroman10-bold.otf,ItalicFont=lmroman10-italic.otf,BoldItalicFont=lmroman10-bolditalic.otf]
\setmonofont{lmmono10-regular.otf}[Scale=MatchLowercase]
\usepackage[margin=25mm]{geometry}
\usepackage[unicode,hidelinks]{hyperref}
\setlength{\emergencystretch}{3em}
\setlength{\parskip}{0.4em}
\setlength{\parindent}{0pt}
\allowdisplaybreaks[2]
\raggedbottom
\title{\texorpdfstring{Lipschitz energy and the averaged lifted potential}{Lipschitz energy and the averaged lifted potential}}
\author{LangevinLean mathematical exposition}
\date{}
\begin{document}\maketitle
\noindent Source module: \texttt{Proximal\allowbreak{}Langevin/\allowbreak{}Splitting\allowbreak{}Coupling.\allowbreak{}lean}. The accompanying HTML pairs this mathematical exposition with the exact Lean statements and records the scope of the formalization.


Let \(g\) be convex, bounded below and \(K\)-Lipschitz, \(\lambda>0\), and \(N>0\). The module proves energy estimates for the true pairwise proximal map and a value-Lipschitz bound for \(G=N^{-1}\sum_i g(J_i z)\).


\section{\texorpdfstring{Obtain energy from value Lipschitzness}{Obtain energy from value Lipschitzness}}


At \(p=P_\lambda(x)\), the residual support at a comparator \(a\) controls \(f(p)-f(a)\). Bound \(f(x)-f(p)\le K\|x-p\|\), expand \(\|x-a\|^2\), and use \((\|x-p\|-\lambda K)^2\ge0\). Thus \(\|p-a\|^2\le\|x-a\|^2-2\lambda(f(x)-f(a))+\lambda^2K^2\). Sum over pairs, use Jensen for the shared parameter, and divide by \(N\) to get the exact averaged energy inequality with remainder \(\lambda^2K^2\).




The pointwise input is a genuine proximal minimizer \(p\). Its residual \(r=(x-p)/\lambda\) is a subgradient at \(p\) and has norm at most \(K\) when the potential is \(K\)-Lipschitz. Combining the support at \(p\), \(f(x)-f(p)\le K\|x-p\|\), and the quadratic expansion yields \[\|p-y\|^2\le\|x-y\|^2-2\lambda(f(x)-f(y))+\lambda^2K^2.\] Sum this for the original pairs and use parameter averaging to pass to the scaled state. The input is value Lipschitzness; no bounded subgradient selection or unproved lifted-proximal identity is hidden in the estimate.




\textbf{Correspondence with the displayed declarations.}


\begin{itemize}

\item \texttt{is\allowbreak{}Prox\allowbreak{}Point\_\allowbreak{}energy\_\allowbreak{}le\_\allowbreak{}of\_\allowbreak{}lipschitz}: The proximal energy inequality follows directly from a Lipschitz potential.
\item \texttt{pairwise\_\allowbreak{}prox\_\allowbreak{}energy\_\allowbreak{}le\_\allowbreak{}lipschitz}: The nonsmooth energy bound for the actual averaged pairwise proximal map, with a global Lipschitz hypothesis instead of an assumed subgradient selection.
\end{itemize}


{\raggedright\noindent Lean declarations: \texttt{is\allowbreak{}Prox\allowbreak{}Point\_\allowbreak{}energy\_\allowbreak{}le\_\allowbreak{}of\_\allowbreak{}lipschitz}, \texttt{pairwise\_\allowbreak{}prox\_\allowbreak{}energy\_\allowbreak{}le\_\allowbreak{}lipschitz}.\par}

\section{\texorpdfstring{Averaged pair distances control the physical-\allowbreak{}time potential}{Averaged pair distances control the physical-time potential}}


Define \(G(z)=F_g(z)/N\). Jensen and \(\sum_i\|J_i v\|^2=N\|v\|^2\) imply \(N^{-1}\sum_i\|J_i v\|\le\|v\|\). Therefore

\[|G(z)-G(w)|\le\frac KN\sum_i\|J_i(z-w)\|\le K\|z-w\|.\]

The same constant survives the averaging because the latent coordinates were scaled by \(N^{-1/2}\).




Let \(\bar f_N(z)=N^{-1}\sum_i f(A_i z)\). Cauchy–Schwarz and the replication identity give \[\frac1N\sum_i\|A_i v\|\le\left(\frac1N\sum_i\|A_i v\|^2\right)^{1/2}=\|v\|.\] Applying \(|f(a)-f(b)|\le K\|a-b\|\) to each pair and summing proves \(|\bar f_N(z)-\bar f_N(w)|\le K\|z-w\|\). The declarations separate the average norm estimate from the value estimate so that the particle scaling can be checked independently.




\textbf{Correspondence with the displayed declarations.}


\begin{itemize}

\item \texttt{averaged\allowbreak{}Potential}: The averaged lifted potential in the physical-time Langevin equation.
\item \texttt{average\_\allowbreak{}joint\_\allowbreak{}norm\_\allowbreak{}le}: Average pair distances are controlled by the scaled state norm.
\item \texttt{averaged\allowbreak{}Potential\_\allowbreak{}lipschitz}: Averaging in the scaled coordinates preserves the pair potential's Lipschitz constant.
\end{itemize}


{\raggedright\noindent Lean declarations: \texttt{averaged\allowbreak{}Potential}, \texttt{average\_\allowbreak{}joint\_\allowbreak{}norm\_\allowbreak{}le}, \texttt{averaged\allowbreak{}Potential\_\allowbreak{}lipschitz}.\par}

\end{document}
